\section{DreamFusion: una aplicación representativa de SDS}

\subsection{Resumen de DreamFusion}

DreamFusion (Poole et al., 2022) es el primer trabajo que aplicó con éxito SDS para la generación 3D a partir de texto. Realiza la transferencia de conocimiento de un modelo difusivo de imagen a texto preentrenado, Imagen, a un NeRF, logrando generar modelos 3D únicamente a partir de descripciones textuales.

\begin{importantbox}{Flujo de trabajo de DreamFusion}
El flujo de trabajo de DreamFusion es el siguiente:
\begin{enumerate}
    \item Inicializar una representación aleatoria de NeRF, $\theta$
    \item En cada iteración:
    \begin{itemize}
        \item Seleccionar aleatoriamente un punto de vista de cámara $\pi$
        \item Renderizar la imagen 2D bajo dicho punto de vista, $x_0 = g(\theta, \pi)$
        \item Calcular el gradiente de SDS y actualizar $\theta$
    \end{itemize}
    \item Tras miles de iteraciones, $\theta$ converge a un modelo 3D de alta calidad
\end{enumerate}
\end{importantbox}

\subsection{Ejemplos de resultados generados}

DreamFusion puede generar modelos 3D basándose en diversas descripciones textuales creativas, tales como:
\begin{itemize}
    \item ``a fox wearing a sweater'' (una zorra usando un suéter)
    \item ``a ghost eating a hamburger'' (un fantasma comiendo una hamburguesa)
    \item ``a DSLR photo of a peacock on a surfboard'' (una fotografía de cámara DSLR de un pavón sobre una tabla de surf)
\end{itemize}

Estos son escenarios creativos para los cuales es casi imposible encontrar modelos 3D correspondientes en internet; sin embargo, dado que los modelos difusivos de imagen se han entrenado con miles de millones de imágenes, han ``visto'' suficientes conceptos visuales relevantes y, por lo tanto, pueden generar resultados 3D coherentes guiados por SDS.

\subsection{Aplicaciones amplias de SDS}

SDS no se limita a la generación 3D; su marco puede extenderse a cualquier tarea que cumpla con las siguientes condiciones:
\begin{enumerate}
    \item Los datos objetivo pueden convertirse en imágenes mediante algún \textbf{mapeo diferenciable}
    \item Existe un \textbf{modelo difusivo preentrenado} disponible
\end{enumerate}

\begin{table}[H]
\centering
\begin{tabular}{lll}
\toprule
\textbf{Aplicación} & \textbf{Tipo de dato} & \textbf{Función de renderizado $g$} \\
\midrule
Texto a 3D & NeRF/Malla 3D & Renderizador 3D diferenciable \\
Generación de gráficos vectoriales & Parámetros de trazos SVG & Rasterización diferenciable \\
Generación 4D & Escenas 3D dinámicas & Renderizador de video \\
Edición 3D & Modelos 3D existentes & Renderizador 3D diferenciable \\
Generación de texturas & Texturas superficiales mapeadas & Mapeo UV y renderizado \\
\bottomrule
\end{tabular}
\caption{Instanciaciones del marco SDS en diversos escenarios de aplicación}
\end{table}

\begin{knowledgebox}{De la difusión de imágenes a la difusión de video}
La idea de SDS también puede combinarse con \textbf{modelos difusivos de video} preentrenados para generar contenido 4D (escenas 3D dinámicas). Al utilizar los fotogramas de video como objetivo de renderizado, se aprovecha el conocimiento sobre la consistencia temporal aprendido por el modelo difusivo de video para generar escenas 3D dinámicas.
\end{knowledgebox}

\subsection{Resumen de la sección}

DreamFusion representa una aplicación semillera de SDS en la generación 3D a partir de texto, demostrando el enorme potencial de transferir grandes previas 2D al dominio 3D. El marco de SDS posee una alta generalidad y puede extenderse a diversas modalidades como gráficos vectoriales, escenarios 4D, texturas y audio.

\newpage
